%==============================================================================
% Phase-Margin Certification of Multi-Agent LLM Systems
% A Behavioral Frequency-Response Approach to Loop-Stability Diagnosis
%
% White paper draft.  Compile with:
%     pdflatex phase_margin_llms
%     bibtex   phase_margin_llms
%     pdflatex phase_margin_llms
%     pdflatex phase_margin_llms
%==============================================================================

\documentclass[11pt,a4paper]{article}

% ---------- Packages ---------------------------------------------------------
\usepackage[margin=1in]{geometry}
\usepackage{amsmath,amssymb,amsthm,amsfonts,mathtools}
\usepackage{bm}
\usepackage{microtype}
\usepackage[T1]{fontenc}
\usepackage{lmodern}
\usepackage{booktabs}
\usepackage{enumitem}
\usepackage{hyperref}
\usepackage[capitalise,noabbrev]{cleveref}
\usepackage[round,authoryear]{natbib}
\usepackage{xcolor}
% (algorithm-like floats below are constructed manually as figures
%  to avoid the texlive-science dependency on algorithm.sty.)

\hypersetup{
  colorlinks=true,
  linkcolor=blue!50!black,
  citecolor=blue!50!black,
  urlcolor=blue!50!black
}

% ---------- Theorem environments --------------------------------------------
\newtheorem{theorem}{Theorem}[section]
\newtheorem{lemma}[theorem]{Lemma}
\newtheorem{proposition}[theorem]{Proposition}
\newtheorem{corollary}[theorem]{Corollary}
\theoremstyle{definition}
\newtheorem{definition}[theorem]{Definition}
\newtheorem{assumption}[theorem]{Assumption}
\newtheorem{example}[theorem]{Example}
\theoremstyle{remark}
\newtheorem{remark}[theorem]{Remark}

% ---------- Notation macros --------------------------------------------------
\newcommand{\R}{\mathbb{R}}
\newcommand{\C}{\mathbb{C}}
\newcommand{\N}{\mathbb{N}}
\newcommand{\Z}{\mathbb{Z}}
\newcommand{\E}{\mathbb{E}}
\newcommand{\Prob}{\mathbb{P}}
\newcommand{\T}{^{\!\top}}
\newcommand{\embspace}{\mathcal{Z}}
\newcommand{\agent}{\mathcal{L}_{\theta}}
\newcommand{\env}{\mathcal{E}}
\newcommand{\obs}{\mathcal{H}}
\newcommand{\sectorial}[2]{\mathcal{S}(#1,#2)}
\newcommand{\probe}{\mathsf{Probe}}
\newcommand{\PhaseMargin}{\Phi_{\mathrm{margin}}}
\newcommand{\Wone}{W_{1}}
\DeclareMathOperator{\diag}{diag}
\DeclareMathOperator{\sgn}{sgn}
\DeclareMathOperator*{\argmin}{arg\,min}

% ---------- Title ------------------------------------------------------------
\title{\textbf{Phase-Margin Certification of Multi-Agent LLM Systems:\\
A Behavioral Frequency-Response Approach\\
to Loop-Stability Diagnosis}}

\author{%
  Mohamed A.~Mabrok\\[2pt]
  \small\texttt{m.a.mabrok@gmail.com}
}

\date{\today}

\begin{document}
\maketitle

% =============================================================================
\begin{abstract}
\noindent
Multi-agent large-language-model (LLM) systems have entered production yet fail
empirically between $41\%$ and $86.7\%$ of the time on standard benchmarks
\citep{Cemri2025MAST}.  A substantial fraction of those failures
are not reasoning errors but \emph{closed-loop dynamical failures}:  agents
oscillate, repeat steps, drift into echo-chamber consensus, or run away from
themselves until their iteration budget is exhausted.
The Multi-Agent System failure Taxonomy (MAST) catalogues fourteen named
failure modes across $1{,}600+$ traces with inter-annotator
$\kappa = 0.88$, and several explicitly correspond to instabilities of the
agent--environment feedback loop.  Yet there exists no pre-flight
\emph{diagnostic} that would predict, before a costly run, whether a given
LLM~+~prompt~+~tool~+~loop will converge, oscillate, or diverge.

We propose such a diagnostic by lifting the recently developed \emph{phase
theory of MIMO LTI systems} \citep{Chen2023PhaseTheory,Wang2024FirstFiveYears} to the
sampled-time, stochastic, behavioural setting required by black-box LLM
agents.  We define a \emph{behavioural phase response} of an LLM agent,
identifiable from a small number of carefully designed probe perturbations in
semantic-embedding space, and prove a \emph{Sampled Behavioural Small-Phase
Theorem} that converts the identified phase response into a quantitative
phase margin.  The margin classifies the agentic loop into one of three
regimes that match the empirical taxonomy of
\citet{Tacheny2025AgenticLoops} and the limit-cycle observations of
\citet{Wang2025AttractorCycles}:  \emph{contractive} (will converge),
\emph{oscillatory} (will cycle), or \emph{exploratory} (will diverge).  The
diagnostic returns a single number from $\mathcal{O}(10^{2})$ rollouts and
addresses approximately six of the fourteen MAST failure categories without
access to model weights, attention activations, or training data.

This paper presents the formalism, the theorem, the
identification protocol, the mapping to documented failure modes, and an
honest catalogue of open problems.  We argue this is the first
control-theoretic instrument that turns the loop-stability fraction of
multi-agent LLM failure from a post-hoc surprise into a measurable margin.
\end{abstract}

\medskip
\noindent\textbf{Keywords.}~Multi-agent LLM systems, phase theory, negative
imaginary systems, behavioural systems, frequency-response identification,
loop stability, agentic AI, robustness certification.

% =============================================================================
\section{Introduction}\label{sec:intro}

The deployment of multi-agent LLM systems---ReAct-style agents, multi-LLM
debate, agentic coding assistants, autonomous research aides---has moved
faster than the diagnostics that should accompany any feedback system at
scale.  This paper is about closing one specific gap in those diagnostics.

\subsection{The empirical problem}

Multi-agent LLM frameworks (AutoGen, MetaGPT, ChatDev, HyperAgent,
AppWorld, AG2, Magentic-One)  fail empirically between $41\%$ and $86.7\%$
of the time on standard benchmarks, depending on the framework and the
task~\citep{Cemri2025MAST}. A peer-reviewed study at NeurIPS~2025 catalogues
$1{,}600+$ annotated execution traces into $14$ named failure modes
clustered in $3$ categories (\emph{system design}, \emph{inter-agent
misalignment}, \emph{task verification}), and validates the taxonomy with
inter-annotator agreement $\kappa = 0.88$.  Two of the named modes,
\texttt{step\_repetition} and \texttt{reasoning\_action\_mismatch},
correspond directly to closed-loop dynamical failures: oscillation between
states and phase decoupling between sub-blocks of the agentic loop.
Independent work has reproduced explicit limit cycles in iterated
paraphrasing~\citep{Wang2025AttractorCycles} and has classified agentic-loop
trajectories into contractive, oscillatory, and exploratory dynamical
regimes~\citep{Tacheny2025AgenticLoops}.  In multi-agent debate
specifically, homogeneous-agent panels collapse into echo-chamber
consensus~\citep{Oh2025BeliefEntrenchment}, a single adversarial agent
can shift system accuracy by $10\%$--$40\%$~\citep{Sci2026PersuasionFailure},
and current debate methods fail to consistently improve over single-agent
baselines~\citep{ICLRBlog2025MAD}.  Production frameworks have responded by
shipping emergency hard-stops such as AutoGen's
\texttt{max\_consecutive\_auto\_reply}, an implicit acknowledgment that the
loop-stability problem is real enough to require an explicit escape hatch.

\subsection{The diagnostic gap}

Despite the documented loss surface, no \emph{pre-flight} diagnostic
predicts loop stability for a given LLM, prompt, and tool configuration.
Practitioners detect failure post-hoc when the iteration cap is reached,
which is the most expensive way to discover divergence.  The folklore
("this prompt is unstable", "this loop tends to thrash") is plausible
but unmeasured.

In every other engineering discipline that puts feedback loops in front of
users---audio, aviation, power grids, automotive---a stability margin is
identified once and read out before deployment.  The instrument behind
that identification is a \emph{frequency-response measurement}: small,
structured perturbations applied at the input, the output is measured,
gain and phase are extracted, and a margin is computed.
The instrument is over a century old in audio~\citep{Bode1945Network}; it
has been generalised to multi-input multi-output systems
\citep{Anderson1973NetworkAnalysis,Megretski1997IQC,Iwasaki2005GeneralizedKYP},
and most recently to phase-bounded systems~\citep{Chen2023PhaseTheory}.
LLM agentic loops are the first widely-deployed feedback system that has
not yet been instrumented with this kind of margin.

\subsection{Contributions}

This paper proposes such an instrument and develops the supporting theory.
Specifically:

\begin{enumerate}[leftmargin=*]
\item We model the LLM agentic loop as a sampled, nonlinear, stochastic
feedback system on a finite-dimensional semantic-embedding manifold
(Section~\ref{sec:agent-model}).
\item We introduce the \emph{behavioural phase response} of an LLM agent:
an empirical phase identifiable from a small number of probe perturbations
of the agent's nominal trajectory in semantic-embedding space, without
access to weights or attention activations
(Section~\ref{sec:behavioural-phase}).
\item We define a corresponding \emph{locally sectorial} class of LLM
agents---the behavioural analogue of the (semi-)sectorial systems of
\citet{Chen2023PhaseTheory}---and prove a \emph{Sampled Behavioural
Small-Phase Theorem} that converts the identified phase response into a
quantitative loop-stability margin
$\PhaseMargin$ (Section~\ref{sec:theorem}).
\item We give a black-box probe protocol that produces $\PhaseMargin$ from
$\mathcal{O}(10^{2})$ rollouts (Section~\ref{sec:protocol}), and predict
the three regimes (contractive / oscillatory / exploratory) reported in
the empirical literature (Section~\ref{sec:regimes}).
\item We map the certificate explicitly onto roughly six of the fourteen
MAST failure modes, and we are explicit about which modes the certificate
does \emph{not} address (Section~\ref{sec:mast-mapping}).
\item We propose an experimental validation plan against the published
benchmarks and the documented limit-cycle datasets
(Section~\ref{sec:validation}), and we list four open theoretical problems
that the framework forces (Section~\ref{sec:open}).
\end{enumerate}

The novelty claim is precise.  The phase theory of LTI systems is in
print~\citep{Chen2023PhaseTheory}; its nonlinear extension is in
print~\citep{Chen2021PhaseNonlinear}; the empirical taxonomy of multi-agent
LLM failure is in print~\citep{Cemri2025MAST}; the geometric description
of agentic loop dynamics is in print~\citep{Tacheny2025AgenticLoops}.  What is
\emph{not} in print, to the best of our literature audit, is a
\emph{sampled-time, stochastic, behavioural} phase certificate identified
from black-box LLM probes that produces a quantitative margin and
matches it to the documented failure modes.  That is what we contribute.

% =============================================================================
\section{The Documented Problem}\label{sec:problem}

We begin by establishing, with citations, that the loop-stability fraction
of multi-agent LLM failures is real, named, measured, and currently
without a pre-flight diagnostic.

\subsection{The MAST taxonomy and quantitative failure rates}

\citet{Cemri2025MAST} construct the first multi-agent system failure
taxonomy, MAST, by annotating $1{,}600+$ execution traces collected from
seven popular multi-agent LLM frameworks across coding, math, and general
agent tasks.  The headline rate is the empirical task-failure rate per
framework, which ranges from $41\%$ to $86.7\%$ on standard benchmarks.
The fourteen failure modes are organised as follows.

\begin{itemize}[leftmargin=*]
\item \emph{System design issues:}  disobeying task specification,
disobeying role specification, step repetition, loss of conversational
history, premature termination.
\item \emph{Inter-agent misalignment:} reasoning-action mismatch,
information withholding, ignored other agent's input, action--reasoning
disagreement, derailing, withholding critical information,
incorrect verification.
\item \emph{Task verification:} insufficient verification,
incorrect verification of action outputs.
\end{itemize}

The two italicised modes \texttt{step\_repetition} and
\texttt{reasoning\_action\_mismatch} correspond directly to feedback-loop
failures.  Together with the qualitative
descriptions in the paper of "conversations cycling without progress",
"perpetual cycle of unproductive exchanges", and "memory limits causing
agents to forget previous discussions" (op.~cit.), they constitute the
loop-stability fraction of the failure surface.

\subsection{Independent reproduction of the dynamical regimes}

Two recent dynamical-systems analyses corroborate the picture.
\citet{Tacheny2025AgenticLoops} formalise the LLM agentic loop as a
discrete-time dynamical system in semantic-embedding space and exhibit
\emph{three classifiable regimes}: contractive trajectories converging to
stable attractors, oscillatory trajectories cycling among attractors, and
exploratory trajectories diverging unboundedly.
\citet{Wang2025AttractorCycles} demonstrate an explicit \emph{2-period
limit cycle} in iterated paraphrasing across multiple frontier LLMs---a
textbook oscillation that would, in any other engineering discipline, be
diagnosed by a phase margin.

\subsection{Multi-agent debate failure}

The picture sharpens for multi-agent debate, the prototypical
multi-agent LLM use case.
\citet{Oh2025BeliefEntrenchment} document that homogeneous-agent
debate panels frequently converge into echo-chamber consensus that
amplifies pre-existing biases instead of correcting them, with the
effect intensifying as panel size grows.  In sequential strategic
environments, a single early suboptimal decision triggers a cascade
that the consensus process cannot self-correct.
\citet{Sci2026PersuasionFailure} quantify adversarial influence:
a single strategically designed adversarial agent reduces system
accuracy by $10\%$--$40\%$ and increases consensus on incorrect
answers by more than $30\%$ on factual benchmarks.
\citet{Wynn2025TalkCheap} catalogue further debate failure modes
including stalemate, premature agreement, and degenerate reasoning.
The synthesis of \citet{ICLRBlog2025MAD} concludes that current
multi-agent debate methods do not reliably outperform single-agent
baselines, even with substantially increased compute---and that simple
majority voting captures most of whatever benefit accrues.

\subsection{Production-side acknowledgement}

The folklore from production deployments matches the academic findings.
The popular AutoGen framework added an explicit
\texttt{max\_consecutive\_auto\_reply} parameter as an emergency hard-stop;
several agent-deployment vendors have published essays titled "the agent
loop problem" or "stop the loop".  The major LLM providers all impose
default iteration caps on their agent products---an implicit confession
that the underlying systems do not converge unaided.

\subsection{The diagnostic gap}\label{ssec:gap}

In every cited reference above the failure modes are detected
\emph{after} the loop has run.  No published method takes an LLM, a
prompt, a tool stack, and a loop topology as input and returns a
quantitative pre-flight stability margin as output.  The
\emph{description} of the regimes is in print; the \emph{prediction
instrument} is not.  Filling this gap is the goal of the present paper.

% =============================================================================
\section{Control-Theoretic Background}\label{sec:background}

We summarise, at the level of working definitions, the four control-theoretic
ingredients we draw on: positive realness and passivity, negative imaginary
systems, the phase theory of \citet{Chen2023PhaseTheory}, and the behavioural
systems framework of \citet{Willems1986Behavioural}.

\subsection{Positive realness and passivity}\label{ssec:pr}

A real-rational proper transfer matrix $G:\C\to\C^{m\times m}$ is
\emph{positive real} (PR) if it is analytic in $\Re s>0$ and
$G(s)+G(s)^{*}\succeq 0$ for $\Re s>0$.  By the
Kalman--Yakubovich--Popov (KYP) lemma~\citep{Anderson1973NetworkAnalysis}
this is equivalent to the existence of $P=P\T \succ 0$ certifying a
state-space dissipation inequality with supply rate
$s(u,y)=u\T y$.  Two PR systems in negative feedback are stable under
mild non-degeneracy, the original "passivity theorem".

\subsection{Negative imaginary systems}\label{ssec:ni}

A real-rational proper $G$ is \emph{negative imaginary}
(NI)~\citep{Lanzon2008NI,Mabrok2014NIFoundations} if it is analytic in
$\Re s>0$ and $j\big(G(j\omega)-G(j\omega)^{*}\big)\succeq 0$ for all
$\omega>0$ where $j\omega$ is not a pole, equivalently
$\Im G(j\omega)\le 0$.  NI captures lightly damped flexible-structure
plants (collocated force--position) for which PR fails.  Two NI systems
in positive feedback are stable under a DC-gain
condition $\lambda_{\max}(M(0)N(0))<1$.

\subsection{Phase theory of MIMO LTI systems}\label{ssec:phase-theory}

\citet{Chen2023PhaseTheory} unify PR and NI into a single phase-bounded
class.  Briefly, a square real-rational $G$ is \emph{(semi-)sectorial} if
$G(j\omega)$ is sectorial as a complex matrix at every frequency, and
its \emph{phase response}
$\phi(G(j\omega))\subset(-\pi,\pi]$ is well-defined and lies in a frequency-dependent
phase sector.  PR corresponds to phase contained in
$[-\pi/2,\pi/2]$; NI to phase in $[-\pi,0]$.  The
\emph{Small Phase Theorem} of~\citet{Chen2023PhaseTheory} states that
$[G_{1},G_{2}]_{-}$ is internally stable if
$\sup_{\omega}\bigl(\overline\phi(G_{1}(j\omega))+\overline\phi(G_{2}(j\omega))\bigr)<\pi$
and a non-encirclement condition holds, where $\overline\phi$ is the upper
phase.  An accompanying \emph{Sectored Real Lemma} provides an LMI
characterisation, and a nonlinear extension exists for input--output
sectorial systems~\citep{Chen2021PhaseNonlinear}.  Mixed gain--phase
robustness is treated in~\citep{Zhao2022SmallGainSmallPhase}.

These results are the foundation we will lift.

\subsection{Behavioural systems theory}\label{ssec:behavioural}

\citet{Willems1986Behavioural,Willems2007Behavioural} pioneered the view
that a dynamical system is defined by the set of admissible
input--output trajectories ("the behaviour"), independent of any
particular state-space realisation.  This is the right standpoint for
LLM agents, whose internal mechanism (weights, attention activations)
is typically inaccessible and irrelevant for closed-loop stability:  what
matters is the input--output map at the embedding level.  The
data-driven realisation of behavioural ideas in
DeePC~\citep{Coulson2019DeePC} and the fundamental lemma of
\citet{Willems2005FundamentalLemma} demonstrate that meaningful
control-theoretic objects can be extracted from data alone.

% =============================================================================
\section{The LLM Agentic Loop as a Sampled Feedback System}\label{sec:agent-model}

We now formalise the object the certificate will measure.

\subsection{State-space in semantic-embedding coordinates}\label{ssec:embedding}

Fix a frozen embedding map $\Pi:\Sigma^{*}\to\embspace\subset\R^{d}$ from
token sequences to a finite-dimensional semantic vector space (e.g.
sentence-T5, the LLM's own pre-LM-head residual stream, or any
fixed-vocabulary representation).  Let $z_{k}\in\embspace$ denote the
embedded running context at iteration $k$ of an agentic loop, and
$y_{k}\in\R^{p}$ a measurable observable extracted from $z_{k}$ via a
linear or smooth nonlinear projection $\obs:\embspace\to\R^{p}$
(e.g. a goal-direction projection, a tool-call type indicator,
or a probe class label).

\begin{assumption}[Embedding regularity]\label{a:embedding}
The embedding map $\Pi$ and the observation map $\obs$ are
$L$-Lipschitz on a neighbourhood of the nominal trajectory
$\{z_{k}^{\star}\}_{k}$ used to compute probes.
\end{assumption}

\subsection{Closed-loop dynamics}\label{ssec:closed-loop}

The agentic loop is the sampled-time discrete dynamical system
\begin{align}
z_{k+1} &= \agent\!\big(z_{k},\,u_{k},\,\xi_{k}\big),
   & y_{k}&=\obs(z_{k}), &
\label{eq:agent}
\\
u_{k+1} &= \env\!\big(y_{k},\,w_{k}\big),
\label{eq:env}
\end{align}
where $\agent:\embspace\times\R^{m}\times\R^{q}\to\embspace$ is the LLM
embedded in the loop topology (decoder + tool router + memory),
$\env:\R^{p}\times\R^{r}\to\R^{m}$ is the environment / tool block
(possibly itself another LLM), $\xi_{k}$ collects the LLM's stochastic
decoding sources, and $w_{k}\sim P_{w}$ is exogenous environmental
noise.

The closed-loop map is
\begin{equation}\label{eq:closed-loop}
z_{k+1} = \Phi(z_{k},z_{k-1},\xi_{k},w_{k})
:= \agent\!\Big(z_{k},\,\env\big(\obs(z_{k-1}),w_{k}\big),\,\xi_{k}\Big).
\end{equation}
Equation~\eqref{eq:closed-loop} is a discrete-time, nonlinear, stochastic
dynamical system on $\embspace\subset\R^{d}$.

\subsection{Why existing phase theory does not directly apply}\label{ssec:why-not}

The phase theory of \citet{Chen2023PhaseTheory} is LTI and continuous-time;
its nonlinear extension~\citep{Chen2021PhaseNonlinear} is
deterministic.  The map $\agent$ in~\eqref{eq:agent} is

\begin{enumerate}[leftmargin=*,label=(\roman*)]
\item \textbf{Discrete-time and sampled.}  The LLM is queried at integer
iterations; the natural frequency variable is $\Omega\in(0,\pi]$ in
samples-per-iteration units.
\item \textbf{Nonlinear and not analytically known.}  $\agent$ is not
expressible in closed form; only black-box rollouts are available.
\item \textbf{Stochastic.}  Decoding produces output variance even with
identical input; $\xi_{k}$ contributes a non-trivial $\sigma_{w}^{2}$.
\item \textbf{Time-varying along the trajectory.}  $\agent$'s effective
linearisation depends on the running context $z_{k}^{\star}$; the
behaviour is at best \emph{locally} Linear Parameter-Varying.
\item \textbf{Defined on a high-dimensional embedding manifold.}  Phase
must be defined directionally with respect to a probing basis in
$\embspace$.
\end{enumerate}

These five features force a behavioural, sampled, stochastic
generalisation of the small-phase theorem.  Sections~\ref{sec:behavioural-phase}
and~\ref{sec:theorem} construct that generalisation.

% =============================================================================
\section{Behavioural Phase Response of an LLM Agent}\label{sec:behavioural-phase}

We define the phase response of $\agent$ not by differentiation of an
analytic model but by observation of input--output behaviour around a
nominal trajectory.

\subsection{Probe protocol input}\label{ssec:probe-input}

Fix a nominal trajectory
$\{(z_{k}^{\star},u_{k}^{\star})\}_{k=0}^{N-1}$ obtained by running the
loop on a representative task.  Pick a unit perturbation direction
$v\in\embspace$, $\|v\|=1$, a probe amplitude $\epsilon>0$, and a probe
discrete frequency $\Omega\in(0,\pi]$.  Inject the perturbation at the
agent's input,
\begin{equation}\label{eq:probe-input}
u_{k}(\Omega,v) = u_{k}^{\star} + \epsilon\,v\cos(\Omega k),\qquad k=0,\dots,N-1,
\end{equation}
and run $M$ independent stochastic rollouts of the open-loop agent
$\agent$ (i.e. with $\env$ replaced by the deterministic playback of
$u_{k}(\Omega,v)$) of length $N$ each.  Let
$\bar y_{k}(\Omega,v) := \frac{1}{M}\sum_{m=1}^{M}\obs(z_{k}^{(m)})$ be
the seed-averaged observable.

\subsection{Empirical phase identification}\label{ssec:phase-id}

Define the residual against the nominal,
$\Delta y_{k}(\Omega,v) := \bar y_{k}(\Omega,v) - \obs(z_{k}^{\star})$, and
fit a sinusoidal model with phase $\theta$, amplitude $A$:
\begin{equation}\label{eq:phase-fit}
\bigl(\widehat\theta_{v}(\Omega),\,\widehat A_{v}(\Omega)\bigr)
:= \argmin_{\theta\in(-\pi,\pi],\,A\ge 0}
   \sum_{k=0}^{N-1}\Big\|\Delta y_{k}(\Omega,v) - \epsilon\,A\cos(\Omega k+\theta)\,b(v)\Big\|_{2}^{2},
\end{equation}
where $b(v):=\partial\obs/\partial z\big|_{z_{k}^{\star}}\,v\in\R^{p}$ is
the probe response direction in observable space (precomputed once via
finite differences if $\obs$ is not given).  Equation~\eqref{eq:phase-fit}
is a complex linear regression, solvable in closed form.

\begin{definition}[Behavioural phase of an LLM agent]\label{def:phase}
The \emph{behavioural phase response} of $\agent$ around
$\{z_{k}^{\star}\}$ in direction $v$ at frequency $\Omega$ is
$\widehat\theta_{v}(\Omega)\in(-\pi,\pi]$ obtained
from~\eqref{eq:phase-fit}.  We similarly define
$\widehat\phi_{v}(\Omega)$ for the environment $\env$ by perturbing
$y_{k-1}$ and observing the induced perturbation in $u_{k}$.
\end{definition}

\subsection{Locally sectorial agents}\label{ssec:locally-sectorial}

\begin{definition}[Locally sectorial agent]\label{def:loc-sectorial}
The agent $\agent$ is \emph{locally sectorial of class}
$\sectorial{\alpha}{\beta}$ around $\{z_{k}^{\star}\}$ in direction $v$ at
fit residual level $\eta>0$ if for every $\Omega\in(0,\pi]$,
\begin{equation}\label{eq:loc-sec}
\widehat\theta_{v}(\Omega)\in[\alpha,\beta],
\end{equation}
and the relative residual
$\rho_{v}(\Omega) := \frac{\sum_{k}\|\Delta y_{k}(\Omega,v)
- \epsilon\widehat A_{v}(\Omega)\cos(\Omega k+\widehat\theta_{v}(\Omega))b(v)\|^{2}}
{\sum_{k}\|\Delta y_{k}(\Omega,v)\|^{2}}\le\eta$.
\end{definition}

\begin{remark}\label{r:reduction-LTI}
If $\agent$ happens to be LTI with transfer matrix $G$ on the linearisation
around $\{z_{k}^{\star}\}$ in direction $v$, then for $\epsilon\to 0$,
$M\to\infty$, $\sigma_{w}\to 0$ we recover
$\widehat\theta_{v}(\Omega)\to\arg(b(v)\T G(e^{j\Omega})v)$, the standard
phase response.  Definition~\ref{def:phase} is therefore a behavioural
generalisation that reduces correctly.
\end{remark}

\begin{remark}\label{r:choice-of-direction}
The phase response in Definition~\ref{def:phase} is direction-dependent
because $\agent$ is a high-dimensional MIMO operator.  In practice we
identify it on a probing basis $\{v_{1},\dots,v_{p}\}$ chosen via the
top-$p$ principal directions of the gradient of $\obs$ along the nominal
trajectory, or via task-relevant semantic axes (e.g. answer-stance, plan
specificity, tool-call category).  See Section~\ref{ssec:basis-choice}.
\end{remark}

\subsection{Sample-complexity of the identification}\label{ssec:sample-complexity}

\begin{lemma}[Phase identification rate]\label{l:rate}
Under Assumption~\ref{a:embedding} and bounded decoding noise
$\E[\|\xi_{k}\|^{2}]\le\sigma_{w}^{2}$, suppose the local linearisation
of $\agent$ in direction $v$ at frequency $\Omega$ has gain
$A_{v}^{\star}(\Omega)\ge A_{\min}>0$ and phase
$\theta_{v}^{\star}(\Omega)$.  Then the empirical phase
$\widehat\theta_{v}(\Omega)$ from~\eqref{eq:phase-fit} satisfies
\[
\E\bigl|\widehat\theta_{v}(\Omega)-\theta_{v}^{\star}(\Omega)\bigr|^{2}
\;\le\; \frac{c\,\sigma_{w}^{2}}{\epsilon^{2}\,A_{\min}^{2}\,M\,N}+\mathcal{O}(\eta),
\]
for an absolute constant $c$, with the second term controlling the
local-linearity bias.
\end{lemma}

The lemma is a standard sinusoid-in-Gaussian-noise least-squares
calculation; the proof is sketched in
Appendix~\ref{app:proofs}.  The headline takeaway is that doubling
the probe amplitude or the sample budget halves the phase
identification variance, and the local-linearity bias is $\mathcal{O}(\eta)$
rather than $\mathcal{O}(\eta^{2})$, so the choice of probe amplitude is a
bias-variance trade.

% =============================================================================
\section{The Sampled Behavioural Small-Phase Theorem}\label{sec:theorem}

We now state the main theorem.  Let $W_{1}$ denote the Wasserstein-1
distance between probability measures on $\embspace$.

\begin{assumption}[Local linearity envelope]\label{a:env}
The agent and environment maps $\agent$ and $\env$ admit local
linearisations along the nominal trajectory whose Jacobian operator
norms in the probing basis are bounded above by $L_{\agent}$ and
$L_{\env}$, with residuals bounded by $\eta_{\agent}$ and $\eta_{\env}$ in
the sense of Definition~\ref{def:loc-sectorial}.
\end{assumption}

\begin{assumption}[Stochastic non-degeneracy]\label{a:stoch}
The decoding noise $\xi_{k}$ has a density bounded above and below on
its support, and the map $\agent$ is irreducible in the sense of having
no absorbing region of zero measure.
\end{assumption}

\begin{theorem}[Sampled Behavioural Small-Phase]\label{thm:main}
Suppose, around the nominal trajectory and in every direction $v$ of a
probing basis $\{v_{1},\dots,v_{p}\}$, the agent and environment satisfy
\begin{equation}\label{eq:sectorial-classes}
\agent\in\sectorial{\alpha_{v}}{\beta_{v}},\qquad
\env\in\sectorial{\gamma_{v}}{\delta_{v}},\qquad v\in\{v_{1},\dots,v_{p}\}.
\end{equation}
Define the directional \emph{phase margin}
\begin{equation}\label{eq:margin}
\PhaseMargin(v) := \min\!\big(\,\pi-(\beta_{v}+\delta_{v}),\;
                                    (\alpha_{v}+\gamma_{v})+\pi\,\big),
\end{equation}
and the global margin
$\PhaseMargin := \min_{v\in\{v_{1},\dots,v_{p}\}}\PhaseMargin(v)$.
Under Assumptions~\ref{a:embedding}--\ref{a:stoch}, if
\begin{equation}\label{eq:main-condition}
\PhaseMargin \;>\; C\!\bigl(\eta_{\agent}+\eta_{\env}\bigr) \;+\;
\frac{C'\sigma_{w}}{\epsilon\,A_{\min}}\frac{1}{\sqrt{MN}},
\end{equation}
for explicit constants $C,C'>0$ depending only on
$L_{\agent},L_{\env},L,d,p$, then the closed-loop process
$\{z_{k}\}$ defined by~\eqref{eq:closed-loop} is

\begin{enumerate}[leftmargin=*,label=(\arabic*)]
\item geometrically ergodic with a unique invariant measure $\pi^{\infty}$,
\item Wasserstein-1 contractive at rate
$\rho \le 1 - \PhaseMargin/\pi + \mathcal{O}(\eta_{\agent}+\eta_{\env}+(MN)^{-1/2})$, in the sense that
$\Wone(\mathcal{P}_{k}\mu,\pi^{\infty})\le \rho^{k}\Wone(\mu,\pi^{\infty})$ for any initial distribution $\mu$,
\item the empirical observable $y_{k}$ converges in distribution to a
unique stationary law and the running mean settles within
$\varepsilon$ of its asymptotic value in $\mathcal{O}\bigl(\frac{1}{\PhaseMargin/\pi}\log\frac{1}{\varepsilon}\bigr)$
iterations.
\end{enumerate}
\end{theorem}

A proof sketch of Theorem~\ref{thm:main} is given in
Appendix~\ref{app:proofs}.  The strategy combines three ingredients:
the deterministic small-phase argument
of~\citep{Chen2023PhaseTheory,Chen2021PhaseNonlinear} bounding the
deterministic part of the closed-loop iteration, a local-linearity
perturbation bound paying $\mathcal{O}(\eta_{\agent}+\eta_{\env})$, and a
discrete-time stochastic Lyapunov / drift argument paying
$\mathcal{O}((MN)^{-1/2})$ for the empirical-margin estimation error.

\subsection{Three regimes from one number}\label{sec:regimes}

The contrapositive of Theorem~\ref{thm:main} produces a regime
classification that matches the empirical literature
\citep{Tacheny2025AgenticLoops,Wang2025AttractorCycles,Cemri2025MAST}.

\begin{corollary}[Regime classification]\label{c:regimes}
With $\PhaseMargin$ defined as in~\eqref{eq:margin}, the agentic loop
exhibits
\begin{itemize}[leftmargin=*]
\item \emph{contractive} dynamics (predicted convergence to a unique
attractor) when $\PhaseMargin>0$ with sufficient confidence;
\item \emph{oscillatory} dynamics (predicted limit cycle, the
\citet{Wang2025AttractorCycles} regime) when $\PhaseMargin\approx 0$ in
expectation but its directional minimum is negative on a positive-measure
set of probe directions;
\item \emph{exploratory} / \emph{divergent} dynamics (predicted
unbounded drift) when
$\alpha_{v}+\gamma_{v}<-\pi$ or $\beta_{v}+\delta_{v}>\pi$ on a
positive-measure direction set.
\end{itemize}
\end{corollary}

The three regimes are exactly the three observed in
\citet{Tacheny2025AgenticLoops}, but here they are predicted from a
quantitative margin computed before the loop is run.

\subsection{Specialisation to known cases}\label{sec:specialisations}

\begin{itemize}[leftmargin=*]
\item \emph{Deterministic LTI agent.}  Setting $\sigma_{w}=0$,
$\eta_{\agent}=\eta_{\env}=0$, $M=N=\infty$, the residual term
in~\eqref{eq:main-condition} vanishes and Theorem~\ref{thm:main} reduces
to the small-phase theorem of~\citet{Chen2023PhaseTheory}.
\item \emph{Deterministic nonlinear agent.}  Setting $\sigma_{w}=0$ but
keeping $\eta>0$, we recover a discrete-time analogue of the nonlinear
small-phase theorem of~\citet{Chen2021PhaseNonlinear}.
\item \emph{Memoryless map} ($\agent$ and $\env$ both phase-zero around
their fixed point).  $\PhaseMargin$ collapses to a small-gain
condition; the iteration is contractive iff $L_{\agent}L_{\env}<1$, which
is the \citet{BanachFixedPoint1922} fixed-point principle.
\end{itemize}

% =============================================================================
\section{The Black-Box Probe Protocol}\label{sec:protocol}

The theorem becomes useful when its hypothesis is checkable from black-box
LLM rollouts.  Algorithm~1 below gives the probe protocol.

\subsection{Algorithm}\label{ssec:algorithm}

\begin{figure}[t]
\centering
\hrule\vspace{4pt}
\noindent\textbf{Algorithm 1: Black-box phase-margin certification of an LLM agentic loop.}
\label{alg:probe}
\vspace{2pt}\hrule\vspace{6pt}
\small\noindent
\textbf{Input:}\, LLM API $\agent$, environment / tool block $\env$,
observable $\obs$, nominal task prompt $p$, probing basis
$V=\{v_{1},\dots,v_{p}\}\subset\embspace$, frequency grid
$\mathcal{G}\subset(0,\pi]$, probe amplitude $\epsilon$, sample budget
$M$, horizon $N$.

\medskip
\begin{enumerate}[label=\textbf{\arabic*.},leftmargin=*,itemsep=2pt,topsep=2pt]
\item Run $M_{0}$ nominal closed-loop rollouts of length $N$ and compute
$\{z_{k}^{\star}\}$ as the running mean trajectory.
\item \textbf{For} each $v\in V$ and each $\Omega\in\mathcal{G}$ \textbf{do}:
\begin{enumerate}[label=\,\textbf{\alph*.},leftmargin=2em,itemsep=2pt,topsep=2pt]
\item Inject $u_{k}(\Omega,v)=u_{k}^{\star}+\epsilon v\cos(\Omega k)$ in
the open-loop agent for $M$ stochastic rollouts of horizon $N$.
\item Compute residuals $\Delta y_{k}(\Omega,v)$ and fit
$(\widehat\theta_{v}(\Omega),\widehat A_{v}(\Omega))$
via~\eqref{eq:phase-fit}.
\item Compute fit quality $\rho_{v}(\Omega)$.  If $\rho_{v}(\Omega)>\eta$
for many $\Omega$, expand the probing basis or reduce $\epsilon$.
\end{enumerate}
\item Repeat steps 2--3 for $\env$ to obtain $\widehat\phi_{v}(\Omega)$.
\item Compute $\alpha_{v},\beta_{v},\gamma_{v},\delta_{v}$ as
$\min/\max$ over $\Omega\in\mathcal{G}$ with statistical buffers from
Lemma~\ref{l:rate}.
\item Compute $\PhaseMargin(v)$ via~\eqref{eq:margin} and
$\PhaseMargin=\min_{v}\PhaseMargin(v)$.
\end{enumerate}

\medskip\noindent
\textbf{Output:}\, $\PhaseMargin$ and the predicted regime
(contractive / oscillatory / exploratory) from
Corollary~\ref{c:regimes}.
\vspace{4pt}\hrule
\end{figure}

\subsection{Probing basis selection}\label{ssec:basis-choice}

The choice of $V=\{v_{1},\dots,v_{p}\}$ trades off coverage against probe
budget.  Three natural defaults:
\begin{enumerate}[leftmargin=*]
\item \emph{Task-direction basis.}  Use a small set of semantic axes
relevant to the task: answer stance ($v_{\text{stance}}$), plan
specificity ($v_{\text{plan}}$), tool-call category
($v_{\text{tool}}$).  Suitable when domain knowledge of the task is
available.
\item \emph{Empirical PCA basis.}  Run nominal rollouts, compute the
PCA of the observed $\{\Delta\obs(z_{k}^{\star})\}_{k}$ residuals against
a slightly perturbed prompt, take the top $p$ components.  Domain-free
default.
\item \emph{Random Hutchinson basis.}  Sample $v\sim\mathrm{Unif}(\mathbb{S}^{d-1})$ and use Hutchinson-style estimators to
upper-bound $\PhaseMargin$ across all directions in expectation.  Useful
for high-dimensional safety-critical applications.
\end{enumerate}

\subsection{Computational cost}\label{ssec:cost}

The probe protocol requires
$|V|\cdot|\mathcal{G}|\cdot M\cdot N$ open-loop agent invocations for
$\agent$, plus the same for $\env$ if it is not deterministic.  With
$|V|=4$, $|\mathcal{G}|=8$, $M=8$, $N=16$, this is $\approx 4{,}000$ LLM
calls---roughly the cost of two failed agent runs in production---to
predict the success of all subsequent runs on the same loop topology.

% =============================================================================
\section{Mapping to Documented MAST Failure Modes}\label{sec:mast-mapping}

We are explicit about which of the fourteen MAST failure modes the
phase-margin certificate addresses.

\begin{table}[h]
\centering
\small
\renewcommand{\arraystretch}{1.15}
\begin{tabular}{@{}p{0.45\textwidth}p{0.5\textwidth}@{}}
\toprule
\textbf{MAST mode} & \textbf{Predicted by $\PhaseMargin$?} \\
\midrule
\texttt{step\_repetition} & \textbf{Yes.}  Negative directional margin
on a stance / plan direction predicts a 2-cycle limit cycle. \\
\texttt{reasoning\_action\_mismatch} & \textbf{Yes.}  Phase mismatch
between sub-block agents shows up as $\beta_{v}+\delta_{v}\to\pi$ on the
relevant direction. \\
\texttt{loss\_of\_conversation\_history} & \textbf{Partially.} Memory
truncation manifests as a discontinuity in the fit
$\widehat\theta_{v}(\Omega)$; flagged but not fully diagnosed. \\
\texttt{premature\_termination} & \textbf{Partially.} Over-contractive
margin ($\PhaseMargin\to\pi$) flags collapse to a too-quick fixed point. \\
\texttt{disobeying\_task\_specification} & \textbf{No.}  Specification
errors are upstream of the loop dynamics. \\
\texttt{disobeying\_role\_specification} & \textbf{No.}  Same. \\
\texttt{information\_withholding} & \textbf{No.}  Inter-agent
communication content; not a loop-stability property. \\
\texttt{ignored\_other\_agent's\_input} & \textbf{Partially.} Shows up
as low directional gain $\widehat A_{v}(\Omega)$, indicating a near-zero
loop element in that direction. \\
\texttt{action\_reasoning\_disagreement} & \textbf{Yes.}  Internal phase
mismatch within $\agent$ between reasoning and acting sub-modules. \\
\texttt{derailing} & \textbf{Partially.}  Exploratory regime is
predicted by $\alpha_{v}+\gamma_{v}<-\pi$. \\
\texttt{withholding\_critical\_information} & \textbf{No.} \\
\texttt{insufficient\_verification} & \textbf{No.}  Verification quality
is a content property. \\
\texttt{incorrect\_verification} & \textbf{No.} \\
\texttt{incorrect\_verification\_of\_action\_outputs} & \textbf{No.} \\
\bottomrule
\end{tabular}
\caption{Mapping the phase-margin certificate to the MAST taxonomy of
\citet{Cemri2025MAST}.  Of the fourteen named modes, the certificate
fully addresses three, partially addresses four, and is silent on
seven.  The seven on which it is silent are content / specification
issues, not loop-dynamics issues.}
\label{tab:mast-mapping}
\end{table}

The certificate cleanly addresses approximately \emph{half} of the
fourteen MAST modes.  This is the loop-stability fraction of the
failure surface, which is precisely the fraction with no current
pre-flight diagnostic.

% =============================================================================
\section{Validation Plan}\label{sec:validation}

The certificate is falsifiable.  We propose three validation experiments
against published datasets and benchmarks.

\subsection{Experiment 1: Three-regime classification}\label{ssec:exp-regime}

\citet{Tacheny2025AgenticLoops} provide a benchmark of agentic-loop
trajectories labelled into the three dynamical regimes (contractive,
oscillatory, exploratory).  We propose:
\begin{enumerate}[leftmargin=*]
\item Run the probe protocol (Algorithm~1) on each
benchmark configuration to obtain $\PhaseMargin$.
\item Apply the regime predictor of Corollary~\ref{c:regimes}.
\item Evaluate the confusion matrix against the ground-truth regimes.
\end{enumerate}
\textit{Headline metric:} predicted-vs-actual regime accuracy, with the
null baseline being the prior class frequencies.  Expected lift if the
theory is correct: substantial (the regimes are deterministic in
$\PhaseMargin$ at low noise).

\subsection{Experiment 2: 2-cycle prediction in successive paraphrasing}\label{ssec:exp-cycle}

\citet{Wang2025AttractorCycles} demonstrate explicit 2-period limit cycles
in iterated paraphrasing across multiple frontier LLMs.  We propose:
\begin{enumerate}[leftmargin=*]
\item Run the probe protocol with the paraphrasing prompt held fixed
and probe directions chosen along stance / lexical-density axes.
\item Predict 2-cycle vs. fixed-point convergence from the directional
sign of $\PhaseMargin(v)$.
\end{enumerate}
\textit{Headline metric:} per-prompt cycle-vs-no-cycle prediction
accuracy.  The published behaviour is well-controlled (specific prompts
exhibit cycles, others do not), so this is a clean falsification target.

\subsection{Experiment 3: MAST step-repetition trace flagging}\label{ssec:exp-mast}

The MAST dataset \citep{Cemri2025MAST} releases $1{,}600+$ annotated traces.
We propose:
\begin{enumerate}[leftmargin=*]
\item For each MAS framework, run the probe protocol on the framework's
default loop topology with a held-out task.
\item Predict per-task whether \texttt{step\_repetition} or
\texttt{reasoning\_action\_mismatch} will trigger.
\item Compare against the MAST annotations on the held-out task.
\end{enumerate}
\textit{Headline metric:} F1 score for predicting these two MAST modes,
with the baseline being the per-framework base rate.

\subsection{Cost and feasibility}\label{ssec:cost-feasibility}

Each probe run costs $\approx 4{,}000$ LLM calls per loop topology, well
within an academic budget for the three frontier APIs.  All three
benchmarks are publicly released with reproducible evaluation scripts.

% =============================================================================
\section{Limitations and Open Problems}\label{sec:open}

We surface the genuine difficulties.  These are not minor.

\paragraph{O1: Embedding-space dependence.}  $\PhaseMargin$ depends on
$\Pi$ and on the probing basis $V$.  An invariant formulation---e.g.
in terms of the information-geometric Fisher metric of the agent's
output distribution---would give a coordinate-free margin.  This is
open.

\paragraph{O2: Local linearity assumption.}  Definition~\ref{def:loc-sectorial}
requires the residual $\rho_{v}(\Omega)$ to be small.  When the agent's
behaviour is sharply nonlinear (e.g. crossing a tool-call decision
boundary), the fit residual exceeds $\eta$ and the certificate
degrades.  A higher-order behavioural phase, e.g. via Volterra-series or
kernel-Hilbert-space identification, would extend the framework but is
open.

\paragraph{O3: Variance reduction.}  At realistic decoding temperatures
$T=0.7$, $\sigma_{w}^{2}$ may eat the margin.  Variance-reduction
designs (antithetic seeds, common random numbers across frequencies, or
deterministic-decode probes followed by stochastic-decode replays) are
needed to make the protocol tractable at high temperature.

\paragraph{O4: Adversarial robustness.}  The certificate is established
around a nominal trajectory.  An adversarial input distribution voids
the certificate.  The robust-IQC version of phase
theory~\citep{Megretski1997IQC,Zhao2022SmallGainSmallPhase} should
extend to give a robust phase margin, but the lift to behavioural
identification is open.

\paragraph{O5: Multi-agent extensions.}  We have stated the theorem for
a two-block loop $[\agent,\env]_{-}$.  Multi-agent debate involves $k>2$
agents in arbitrary topologies.  The cyclic small-phase
theorem~\citep{Wang2023CyclicSmallPhase} treats the LTI case; its
behavioural-stochastic generalisation is open.

% =============================================================================
\section{Related Work}\label{sec:related}

\textbf{Phase theory of MIMO LTI systems.}  The substrate
of~\citep{Chen2023PhaseTheory,Wang2024FirstFiveYears,Chen2021PhaseNonlinear,Zhao2022SmallGainSmallPhase,Wang2022DiscretePhase,Wang2023LTPPhase,Wang2023CyclicSmallPhase}
is what we lift.  Their framework is LTI (or deterministic nonlinear in
the input-output sectorial sense); our contribution is the sampled,
stochastic, behavioural extension specific to LLM agents.

\textbf{Negative imaginary systems.}  PR/NI/passivity foundations:
\citep{Anderson1973NetworkAnalysis,Lanzon2008NI,Mabrok2014NIFoundations,Patra2011MixedNIBR,Griggs2007MixedSGP}.
The recent neural-ODE application~\citep{Mabrok2025NINeuralODE} controls
mechanical plants with NN controllers carrying the NI property; our
target is loop dynamics of LLM agents, a different object.

\textbf{Behavioural and data-driven systems.} \citet{Willems1986Behavioural,Willems2007Behavioural,Willems2005FundamentalLemma,Coulson2019DeePC}.  These motivate the
black-box-probe formulation of phase identification.

\textbf{Lipschitz / SDP certification of neural networks.}
\citep{Fazlyab2019Lipschitz,Hashemi2024LipShiFT,Yang2024LipNeXt}.  These
methods constrain magnitude (gain), not phase, and treat the network
statically rather than as a closed-loop system.

\textbf{Dynamical-systems analyses of LLM loops.}
\citep{Tacheny2025AgenticLoops,Wang2025AttractorCycles}.  These describe the
regimes; we predict them.

\textbf{Multi-agent LLM failure analyses.}
\citep{Cemri2025MAST,Oh2025BeliefEntrenchment,Sci2026PersuasionFailure,Wynn2025TalkCheap,DuMultiAgentDebate2024,ICLRBlog2025MAD}.
These document and explain the failures; we add a pre-flight predictor.

\textbf{Stable architectures for deep
learning.}~\citep{Haber2017StableArch} construct continuous-time stable
ResNets; orthogonal to the loop-stability question we address.

\textbf{Score-based generative model
stability.}~\citep{ChenSampling2023,Forgetting2026}.  Dissipativity for
diffusion sampling; orthogonal to LLM agent loops.

% =============================================================================
\section{Conclusion}\label{sec:conclusion}

Multi-agent LLM systems have a documented loop-stability problem.
Multi-agent system failure rates of $41\%$--$86.7\%$, a peer-reviewed
$14$-mode taxonomy, independent reproductions of explicit limit cycles,
adversarial-influence accuracy losses of $10$--$40\%$, and emergency
hard-stops shipped in production frameworks all point at the same
conclusion: the closed-loop dynamics of these systems are unstable in
real ways and there is no pre-flight diagnostic.

This paper proposes such a diagnostic by lifting the recently developed
phase theory of MIMO LTI systems to the sampled-time, stochastic,
behavioural setting required by black-box LLM agents.  The Sampled
Behavioural Small-Phase Theorem converts a small number of probe
rollouts into a single number $\PhaseMargin$ that predicts which of
three regimes the loop will exhibit, and the regime classification
matches the empirical taxonomy in the literature.

The certificate addresses approximately half of the fourteen documented
MAST failure modes---specifically the loop-stability half, which has no
current pre-flight diagnostic.  The other half (specification,
verification, content) belongs to other categories of analysis and
remains out of scope.  Five open problems---embedding invariance, local
linearity, variance reduction, adversarial robustness, and the multi-agent
extension---are honestly catalogued.

The bigger picture: every other engineering discipline that puts feedback
systems in front of users has long since had a way to measure stability
margins before the user is exposed.  Bringing that instrument over to
the LLM agent's bench is the small, specific, useful, and---to the best
of our literature audit---genuinely unaddressed contribution.

% =============================================================================
\appendix
\section{Proof sketches}\label{app:proofs}

\subsection*{Lemma~\ref{l:rate}}
The least-squares fit~\eqref{eq:phase-fit} is equivalent to a complex
linear regression on the discrete Fourier basis at frequency $\Omega$.
Stack the residuals $\Delta y_{k}$ along $b(v)$ into a complex
sufficient statistic $S_{N}(\Omega) := \sum_{k=0}^{N-1}\Delta y_{k}\T b(v)\,
e^{-j\Omega k}$.  Under decoding noise of variance $\sigma_{w}^{2}$ on
$\bar y_{k}$ (after seed-averaging by $M$), the noise in $S_{N}$ is
complex Gaussian with variance $N\sigma_{w}^{2}\|b(v)\|^{2}/M$.  The
deterministic part is
$\frac{\epsilon N}{2}A_{v}^{\star}(\Omega)e^{j\theta_{v}^{\star}(\Omega)}\|b(v)\|^{2}+\mathcal{O}(\eta\epsilon N)$.
Phase identification reduces to $\arg(S_{N})$, whose variance is
controlled by the noise-to-signal ratio in the standard
$|\Delta\arg|^{2}\le |\text{noise}|^{2}/|\text{signal}|^{2}$ inequality
\citep{Rife1974FrequencyEstimation}, giving the stated rate. The
$\mathcal{O}(\eta)$ bias term is the local-linearity residual mass leaking
into the wrong frequency bin.\hfill$\square$

\subsection*{Theorem~\ref{thm:main}}
The strategy mirrors the proof of the Small Phase Theorem
in~\citep{Chen2023PhaseTheory} but lifts each of three pieces.

\textbf{(i) Deterministic sectorial composition.}
Around the nominal trajectory in any direction $v$, the
linearised iteration on the residual is
$\widetilde z_{k+1}=J_{v}\widetilde z_{k}$, where $J_{v}$ is the
composition of two locally-linear maps with phase responses in
$[\alpha_{v},\beta_{v}]$ and $[\gamma_{v},\delta_{v}]$.  The composed
phase is in $[\alpha_{v}+\gamma_{v},\beta_{v}+\delta_{v}]$.  By the
deterministic small-phase theorem, the spectral radius
$\rho(J_{v})\le\cos(\PhaseMargin(v)/2)<1$ if
$\PhaseMargin(v)>0$.  This gives the deterministic contraction.

\textbf{(ii) Local-linearity perturbation.}  The residual term
$\eta_{\agent}+\eta_{\env}$ enters as a Lipschitz perturbation of the
deterministic Jacobian; standard perturbation analysis adds
$O(\eta_{\agent}+\eta_{\env})$ to the spectral radius.

\textbf{(iii) Stochastic Lyapunov / drift.}  Under
Assumption~\ref{a:stoch}, the noise injects a bounded random walk into
the iteration.  A standard Foster--Lyapunov drift argument
\citep{MeynTweedie} on $V(z) := \|z-z^{\star}\|^{2}$ gives, for any
$\rho<1-\PhaseMargin/\pi+\mathcal{O}(\eta+(MN)^{-1/2})$,
$\E[V(z_{k+1})\,|\,z_{k}]\le \rho^{2}V(z_{k})+b$ for some
$b<\infty$.  Geometric ergodicity, Wasserstein contraction, and the
mixing-time estimate follow.\hfill$\square$

\section{Discussion of probing-basis sensitivity}\label{app:basis}

A probe basis $V$ that misses a destabilising direction can produce a
falsely positive $\PhaseMargin$.  The Hutchinson random-direction
basis (Section~\ref{ssec:basis-choice}) has the advantage that
$\E_{v\sim\mathrm{Unif}(\mathbb{S}^{d-1})}[\PhaseMargin(v)]$ converges to a
direction-averaged margin under enough samples; the variance bound is
$\mathcal{O}(d/p)$.  Hybrid bases combining task directions and random
directions are recommended for safety-critical use.

\bibliographystyle{plainnat}
\bibliography{references}

\end{document}
direction-averaged margin under enough samples